% vim: spelllang=cs spell textwidth=120
\documentclass{beamer}
\definecolor{links}{HTML}{2A1B81}
\usepackage[utf8]{inputenc}
%\usepackage{ucs}
%\usepackage[utf8x]{inputenc}
\usepackage{minted}
\usepackage{hyperref}
\usepackage[overlay]{textpos}
\hypersetup{colorlinks,urlcolor=links}
\setbeamertemplate{navigation symbols}{}

\begin{document}

\title{Deska: Tool for Central Administration of a Grid Site}
\author{Jan Kundrát \\ \href{mailto:kundratj@fzu.cz}{\nolinkurl{kundratj@fzu.cz}}
\vspace{5mm}
\\ Lukáš Kerpl \\ Martina Krejčová \\ Tomáš Hubík}
\institute{Institute of Physics of the AS CR, v.v.i.}
\begin{frame}
\begin{columns}
\begin{column}{.4\textwidth}
\includegraphics[width=4cm]{../HEPIX-2011/logo-FZU-male.pdf} \\
\includegraphics[width=4cm]{../HEPIX-2011/deska_logo.pdf}
\end{column}
\begin{column}{.6\textwidth}
\maketitle
\end{column}
\end{columns}
\end{frame}


\begin{frame}[fragile]
\frametitle{Úvod}
\begin{itemize}
    \item Popis problému
    \item Navrhované řešení
    \item Databázové objekty
    \item Technická implementace
    \item Problémy
\end{itemize}
\end{frame}


\begin{frame}[fragile]
\frametitle{Problém: duplicitní informace}
\begin{itemize}
    \item Informace o každém serveru / službě rozprostřena po celé infrastruktuře
        \begin{itemize}
            \item HW DB
            \item síťové přepínače
            \item DHCP
            \item DNS
            \item specifikace rolí v Cfengine
            \item Torque a CPU multiplikátory
            \item MRTG \& RRD grafování
            \item Nagios
            \item Ganglia
            \item Munin
            \item \ldots
        \end{itemize}
    \item Komplikované změny, časté omyly, zdroj komplikací
\end{itemize}
\end{frame}


\begin{frame}[fragile]
\frametitle{Navrhované řešení}
\begin{columns}
\begin{column}{.45\paperwidth}
    \begin{itemize}
        \item Jedno centrální místo jako autoritativní zdroj informací
        \item Provádění změn:
            \begin{enumerate}
                \item Změna dat v DB
                \item Vygenerování výstupu
                \item Kontrola změn
                \item Uložení k distribuci
                \item Nasazení do praxe
            \end{enumerate}
        \item Neschválené změny neovlivňují další uživatele ani stav produkčních systémů
    \end{itemize}
\end{column}
\begin{column}{.55\paperwidth}
\begin{textblock}{50}(0,-7.2)
\includegraphics[height=90mm, trim=28mm 53mm 30mm 28mm, clip=true]{../../technical/img-deska-workflow.pdf}
\end{textblock}
\end{column}
\end{columns}
\end{frame}


\begin{frame}[fragile]
\frametitle{Požadavky zadavatele}
\begin{itemize}
    \item CLI rozhraní
    \item WWW přístup není požadován
    \item Skriptování
    \item Verzování, audit logging
    \item Preference FLOSS, Linux
    \item Integrace s existující infrastrukturou
    \item Nevynalézat znovu kolo
\end{itemize}
\end{frame}


\begin{frame}[fragile]
\frametitle{Struktura databáze}
\begin{itemize}
    \item Dynamické požadavky, každé centrum trochu jiné
        \begin{itemize}
            \item Generická databáze
            \item Schéma dotváří správce aplikace
        \end{itemize}
    \item Implementace
        \begin{itemize}
            \item SQL DDL jako definice schématu
            \item Generátory kódu pro uložené procedury
        \end{itemize}
    \item Zapouzdření
        \begin{itemize}
            \item Klientské aplikace {\em nezávislé} na schématu
            \item Komunikační protokol schopný popsat uživatelské schéma
            \item Introspekce
        \end{itemize}
\end{itemize}
\end{frame}


\begin{frame}[fragile]
\frametitle{Databázová struktura: objekty}
\begin{itemize}
    \item Každá položka reprezentovaná {\bf objektem} patřičného {\bf typu} (``kind'')
        \begin{itemize}
            \item Kind $\rightarrow$ databázová tabulka
                \begin{itemize}
                    \item Určuje povolené atributy a jejich datové typy
                    \item {\tt host}, {\tt hardware}, {\tt vendor},\ldots
                \end{itemize}
            \item Objekt $\rightarrow$ řádka v tabulce
                \begin{itemize}
                    \item ({\tt wn123}, {\tt www01}, {\tt wn123->eth0},\ldots)
                \end{itemize}
        \end{itemize}
        \begin{minted}{text}
            host wn123
                hardware sgi-xe310
                rack R06
                position 10
            end
        \end{minted}
    \item Schéma je uživatelsky definováno
        \begin{itemize}
            \item Deska server neobsahuje nic specifického pro datacentrum
        \end{itemize}
\end{itemize}
\end{frame}


\begin{frame}[fragile]
\frametitle{Databázová struktura: relace}
\begin{itemize}
    \item Relace určují pravidla pro vztahy mezi jednotlivými tabulkami
    \item Přidání konkrétního {\em významu} cizím klíčům
        \begin{itemize}
            \item {\bf Template} určuje výchozí hodnotu atributů
                \begin{itemize}
                    \item Dodávka mnoha serverů najednou
                \end{itemize}
            \item Objektová {\bf kompozice}
                \begin{itemize}
                    \item Opětovné využití subkomponent -- {\tt server} i {\tt switch} jsou v {\tt rack}u
                \end{itemize}
            \item {\bf Vnořování} objektů
                \begin{itemize}
                    \item Umožňuje vytvářet hierarchické stromy objektů modelující realitu
                    \item Síťová rozhraní: {\tt srv123->eth0}
                \end{itemize}
            \item {\bf Odkazování}
                \begin{itemize}
                    \item Odkaz na patřičný {\tt warranty\_contract} z každého {\tt hardware}
                \end{itemize}
            \item {\bf Tagování}
                \begin{itemize}
                    \item Přiřazení logických rolí jednotlivým serverům
                \end{itemize}
        \end{itemize}
\end{itemize}
\end{frame}


\begin{frame}[fragile]
\frametitle{Technická implementace}
\begin{columns}
\begin{column}{.45\paperwidth}
    \begin{itemize}
        \item Klient-server aplikace
            \begin{itemize}
                \item 100\% oddělení
            \end{itemize}
        \item Klient:
            \begin{itemize}
                \item CLI
                \item C++, Boost, Spirit
                \item komunikace přes JSON, SSH
            \end{itemize}
        \item Server:
            \begin{itemize}
                \item PostgreSQL 9.0
                \item Pl/PgSQL, PgPython
                \item generátory kódu
            \end{itemize}
        \item Výstupní generátory:
            \begin{itemize}
                \item Python
                \item nativní přístup k DB objektům
                \item Git
            \end{itemize}
    \end{itemize}
\end{column}
\begin{column}{.55\paperwidth}
\begin{textblock}{50}(0,-7.5)
\includegraphics[width=65mm, trim=28mm 64mm 28mm 27mm, clip=true]{../../technical/img-deska-components.pdf}
\end{textblock}
\end{column}
\end{columns}
\end{frame}


\begin{frame}[fragile]
\frametitle{Skriptování v Pythonu}
\begin{itemize}
    \item Zpřístupnění obsahu databáze přes nativní objekty
        \begin{itemize}
            \item Vytvoření hierarchie tříd za běhu
            \item Introspekce přes oficiální komunikační protokol
            \item Čistě automatický proces
        \end{itemize}
    \item Filtrování probíhá na straně serveru
    \item Syntaxe podobná standardnímu ORMu SQL Alchemy
\end{itemize}
\begin{minted}{python}
import deska
# introspekce
deska.init()

for host in deska.host[(deska.host.color == "red") &
                       deska.service.contains("www")]:
    print "%s\n" % host
\end{minted}
\end{frame}


\begin{frame}[fragile]
\frametitle{Problémy v použitých komponentách}
\begin{itemize}
    \item PostgreSQL: nedostatečná chybová diagnostika u selhávajících integritních omezení
        \begin{itemize}
            \item náš patch zařazen do PostgreSQL 9.2
            \item nabídka zaměstnání v evropské společnosti
        \end{itemize}
    \item Boost.Process: špatné ošetření návratových hodnot syscallů {\tt read(2)}, {\tt write(2)}
        \begin{itemize}
            \item neoficiální knihovna, není součástí Boostu
            \item opraveno lokálně, další rozšíření funkčnosti
        \end{itemize}
    \item JSON-Spirit: serializace floating-point čísel neodpovídá specifikaci JSONu
        \begin{itemize}
            \item patch
        \end{itemize}
    \item PsycoPg2: porušení zpětné kompatibility mezi verzemi 2.4.2 a 2.4.1
        \begin{itemize}
            \item workaround, podporujeme všechny rozšířené verze
        \end{itemize}
    \item PgPython: nedostatečné rozlišování prázdných polí a prázdných řetězců, nedostatečná kontrola chyb
        \begin{itemize}
            \item workaround
        \end{itemize}
    \item CMake/CDash: integrace s Redmine
        \begin{itemize}
            \item patch
        \end{itemize}
\end{itemize}
\end{frame}


\begin{frame}[fragile]
\frametitle{Shrnutí}
\begin{itemize}
    \item FIXME
    \item \ldots
\end{itemize}
\end{frame}


\begin{frame}[fragile]
\frametitle{Dotazy}
\begin{center}
    Děkujeme za pozornost!
\end{center}
\end{frame}

\end{document}
